% BTP-I joint report: IEEE conference format (IEEEtran), 6-8 pages.
%
% A GENERATED FIRST DRAFT. Each author rewrites the prose of their own part in
% their own words until Turnitin's AI score is below 20%; the numbers, tables,
% equations and references stay. Status 22 Nov: V and VI rewritten by Anubhav
% (a note above each rewritten subsection lists the edits and added facts).
% 23 Nov: merged with Biraj's version (VII, VIII, hospital details in III and the
% acknowledgment, abstract and conclusion edits, Figs. 4, 5 and 7). Changes made
% while merging, to keep the report consistent and within 8 pages:
%  - V-B: removed the redness-index sentences (no experiment in the report backs
%    them) and the attention-gate description (its channel width is not the one
%    in the extraction code, which uses half the skip channels).
%  - VI-C: removed the hospital-transfer paragraph (kept out of this report;
%    "five-fold cross-validation" was also not the protocol used).
%  - VII-A: removed pixel accuracy and specificity (V-C explains why accuracy is
%    not used for masks).
%  - VII-D: the predicted-vs-measured and residual plots repeat Fig. 6 (a) and
%    (b), so only the absolute-error histogram is kept; text points to Fig. 6.
%  - III: removed the EXIF sentence (VII-A already says it).
%  Later on 23 Nov Biraj sent reworded VII-VIII; that text replaced the earlier
%  one, with the same removals re-applied (listed at the start of VII).
%  - Figure widths reduced: Figs. 2, 3, 4, 5 and 7. Lines starting with % are notes and never
% appear in the PDF. Red [TODO] marks are facts only the authors can supply.
%
% Parts:  shared            title, abstract, I-IV, IX, X, acknowledgment
%         extraction        V, VI
%         Hb model and app  VII, VIII
%
% Overleaf: upload main.tex, refs.bib and the figures listed in README.md into a
% folder called figures/. A figure that is not uploaded yet compiles as a box
% naming the file, so the document always builds.
%
% PAGE BUDGET (8 pages, references included; each section's budget includes its
% own tables and figures):
%   title + abstract 0.3 | I 0.4 | II 0.3 | III 0.25 | IV 0.25   front    1.5
%   V 1.0 | VI 1.5                                                extract  2.5
%   VII 1.5 | VIII 1.5                                            Hb + app 3.0
%   IX 0.35 | X 0.15                                              closing  0.5
%   references                                                             0.5

\documentclass[conference]{IEEEtran}
\usepackage{cite}
\usepackage{amsmath,amssymb}
\usepackage{graphicx}
\usepackage{booktabs}
\usepackage{xcolor}
\usepackage{url}

\newcommand{\todo}[1]{\textcolor{red}{[TODO: #1]}}
% \fig{width}{file}: the figure if present (figures/ on Overleaf, ../figures/
% when compiled inside the repository), otherwise a labelled box.
\newcommand{\fig}[2]{%
  \IfFileExists{figures/#2}{\includegraphics[width=#1]{figures/#2}}{%
  \IfFileExists{../figures/#2}{\includegraphics[width=#1]{../figures/#2}}{%
  \fbox{\parbox[c][2.6cm][c]{\dimexpr#1-2\fboxsep-2\fboxrule\relax}{\centering
  \footnotesize Figure to add:\\ \texttt{\detokenize{#2}}}}}}}

\begin{document}

\title{Anaemia Screening from Smartphone Photographs\\ of the Palpebral Conjunctiva}

\author{
\IEEEauthorblockN{Anubhav Jindal}
\IEEEauthorblockA{Department of Computer Science and Engineering\\ Indian Institute of
Technology Guwahati, Guwahati, India\\ j.anubhav@iitg.ac.in}
\and
\IEEEauthorblockN{Biraj Sarkar}
\IEEEauthorblockA{Department of Computer Science and Engineering\\ Indian Institute of
Technology Guwahati, Guwahati, India\\ s.biraj@iitg.ac.in}
}

\maketitle
% No figure or table at the top of the first page (IEEE practice).
\suppressfloats[t]

% ============================================================ ABSTRACT (shared)
\begin{abstract}
% Assembled by Claude (23 Nov) from Anubhav's own sentences in Sections I and
% VI (the anaemia definition, "We use a phone photo...", "much better than the
% rule-based ones", "too small a range to decide a winner", "can make a
% difference look real when it is not", "the white patch gives ... the best"),
% with the numbers added. The hospital and Hb sentence is condensed from the
% earlier abstract and Section VII (AI-written).
Anaemia is the condition where the haemoglobin (Hb) level in the blood drops
below a certain threshold, and confirming it needs a blood sample. We use a
phone photo to assign a number to the redness of the palpebral conjunctiva, the
inner lining of the lower eyelid. In an independent study on two public
datasets, learned extractors were much better than rule-based ones, with a Dice
score of 0.83 against 0.48 on close-ups of people not seen in training. The top
three of six architectures are within 0.007 Dice of each other, too small a
range to decide a winner, and a single split can make a difference look real
when it is not. Against a white plastic edge present in every photograph, the
white patch gives the best estimate of the light (median 1.5$^\circ$ and
0.8$^\circ$, against about 5$^\circ$ for gray-world). The segmenter selected on
the basis of this study, an Attention U-Net, reaches a Dice of 0.83 on 43
held-out hospital patients, and a DenseNet-121 estimates Hb with a mean
absolute error of 1.35~g/dL on 41 patients, against 1.37~g/dL for predicting
the training mean. A mobile application returns the estimate with a screening
flag.
\end{abstract}

\begin{IEEEkeywords}
anaemia, palpebral conjunctiva, semantic segmentation, U-Net, colour
constancy, haemoglobin estimation, mobile health
\end{IEEEkeywords}

% ============================================================ I. INTRODUCTION (shared)
% NFHS-5 national figures: check 57.0% (women 15-49) and 67.1% (children
% 6-59 months) against the fact sheet before submitting.
\section{Introduction}
% Written by Anubhav (23 Nov, second version). Edits: capitals, "the World
% Health Organization", "threshold ... is below 13.0" -> "is 13.0" (anaemia is
% below it), "introduce ... to" -> "expose ... to", "pulling the eyelid down",
% "causing it to become significantly paler" -> "so it turns noticeably paler
% when Hb is low" (the vessels are why its colour follows Hb), "Few issues" ->
% "There are a few issues", citations. Wording taken from Claude's fact sheet:
% "confirming it needs a blood sample ... is hard" and the three issue bullets.
% Contributions: Anubhav's wording from his first version, shortened by Claude
% to one sentence each. The section-map sentence at the end is Claude's.
Anaemia is the condition where the haemoglobin (Hb) level in the blood drops
below a certain threshold. This threshold, as defined by the World Health
Organization (WHO), is 13.0~g/dL for men, 12.0~g/dL for non-pregnant women, and
11.0~g/dL for pregnant women and children aged 6--59 months~\cite{who2011}. In
India specifically, the fifth National Family Health Survey (NFHS-5) found
anaemia in 57.0\% of women aged 15--49 and 67.1\% of children aged 6--59
months~\cite{nfhs5}. But confirming it needs a blood sample and a lab or a
haemoglobin meter, so screening outside clinics is hard. Also, taking blood can
expose health workers to blood-borne diseases.

Doctors have been using the inner lining of the lower eyelid, called the
palpebral conjunctiva, by pulling the eyelid down slightly and checking its
paleness to predict the presence of anaemia. This tissue has no skin over it
and it is rich in blood vessels, so it turns noticeably paler when Hb is low.
But this judgement is subjective. We use a phone photo to assign a number to
the redness of this tissue. There are a few issues:
\begin{itemize}
  \item the tissue has to be found precisely, because lid skin, sclera or glare
    inside it changes the colour;
  \item the colour depends on the light and on the camera's processing;
  \item labelled data is small, and repeat photos of the same person can end up
    in both training and testing.
\end{itemize}

The main contributions of this work are as follows:
\begin{itemize}
  \item We conduct an independent study of conjunctiva extraction, comparing a
    rule-based approach with learned extraction methods and evaluating six
    segmentation architectures on individuals not included in the training
    data, which guides our choice of segmentation model.
  \item We examine the evaluation methodology through repeated train-test
    splits that keep each individual on one side of every split, a corrected
    significance test and the $F_{0.5}$ score, which reveal limitations in
    conclusions drawn from a single split.
  \item We compare different approaches for estimating illumination in an image
    against a white reference present in every photograph.
  \item We introduce a hospital dataset of 288 patients at ESIC Medical College
    and Hospital, Beltola, with manually annotated conjunctiva masks and
    haemoglobin measurements, used to train the deployed segmentation model and
    a DenseNet-121 model for haemoglobin estimation.
  \item We develop a mobile application, EyeCheck, and a supporting server that
    return an estimated haemoglobin level and a screening flag whose threshold
    takes age, sex and pregnancy status into account.
\end{itemize}
Section~II reviews related work and Section~III the data. Section~IV outlines
the system, Sections~V--VI cover extraction, Section~VII Hb estimation and
Section~VIII the application. Section~IX discusses limitations.

% ============================================================ II. RELATED WORK (shared)
\section{Related Work}
% Written by Claude at Anubhav's request (23 Nov), shorter than the draft it
% replaces. Counts as AI-written text. Sources: papers/conjunctiva-anemia/01-07.
\subsection{Haemoglobin from Conjunctival Photographs}
Collings et~al.~\cite{collings2016} related the redness of the conjunctiva in
digital photographs to Hb, with a colour card in every frame. Suner
et~al.~\cite{suner2021} estimated Hb from colour features of a selected region
in smartphone photographs, and Zhao et~al.~\cite{zhao2024} turned this into a
real-time application. Kato et~al.~\cite{kato2024} applied machine and deep
learning to conjunctival photographs, mostly of children. Mannino
et~al.~\cite{mannino2018} used fingernail photographs and showed that
calibrating to each person improves accuracy. Most of this work controls the
light with a card or raw capture, or selects the region by hand.

\subsection{Segmenting the Conjunctiva}
Kasiviswanathan et~al.~\cite{kasiviswanathan2020} trained segmentation networks
to find the conjunctiva. Dhalla et~al.~\cite{dhalla2023} compared U-Net,
UNet++, FPN, PSPNet and LinkNet on a DenseNet-121 encoder, using 324
photographs of 162 children, and chose LinkNet; their three best models were
within 0.11 percentage points of Dice on a single split, and augmentation before
the split may have put copies of test images into training. We test whether
such rankings hold on repeated, person-safe splits.
person-safe splits.

% ============================================================ III. DATA (shared)
\section{Data}
\begin{table}[t]
\caption{Datasets}
\label{tab:data}
\centering
\footnotesize
\begin{tabular}{@{}lcccc@{}}
\toprule
Dataset & Subjects & Images & Masks used & Hb \\
\midrule
Eyes-Defy-Anemia \cite{dimauro2023} & 218 & 218 & 101 & 217 \\
Mendeley \cite{rahman_mendeley} & $\le$199 & 547 & 547 & -- \\
Hospital (this work) & 288 & 288 & 288 & 279 \\
\bottomrule
\end{tabular}
\end{table}
% Assembled by Claude (23 Nov). Anubhav's own sentences from V-C, V-D and VI-A:
% the device with a white plastic edge, white balance fixed and no flash, "many
% images of the same person's eyes, taken from different angles", full-face
% photographs. Counts and licences are facts. The hospital paragraph is
% Biraj's text (AI-written by his account), shortened. The split sentence is
% from his revised VII-A (23 Nov), verbatim.
Table~\ref{tab:data} lists the three datasets. The
\emph{Eyes-Defy-Anemia}~\cite{dimauro2023} pictures were taken with a device
that has a white plastic edge below the eye; the camera's white balance was
fixed and no flash was used. It has one photograph each of 218 people from
India and Italy, with an Hb value for 217. Many of its annotations are
rectangles, so we used only the 101 traced masks. The Italian photographs come
from two capture batches that differ in age, Hb and light.
\emph{Mendeley}~\cite{rahman_mendeley} has many images of the same person's
eyes, taken from different angles, and many are full-face photographs. Each of
its 547 images has two annotators and no Hb, and a duplicate search puts them at
no more than 199 people. Both public sets are openly licensed (CC~BY-SA~4.0 and
CC~BY~4.0).

The \emph{hospital} cohort has one photograph and a recorded Hb for each of 288
patients at ESIC Medical College and Hospital, Beltola, Guwahati, taken between
September and November 2026 on Xiaomi and Motorola smartphones without flash,
under informed patient consent and the clinical oversight of Dr.~Sandip Paul,
with Hb from the hospital's automated haematology analyser. Every conjunctiva
was outlined by hand. Nine records had no valid Hb, leaving 279 (mean
11.4~g/dL; 106, or 38\%, below 11.0~g/dL). Patients were split by patient registration number into
training (202), validation (43) and test (43) sets, so no patient appears in
more than one set.

% ============================================================ IV. SYSTEM OVERVIEW (shared)
\section{System Overview}
\begin{figure*}[t]
\centering
\fig{0.9\textwidth}{pipeline.png}
\caption{The deployed pipeline. The application sends a photograph to the
server, which segments the conjunctiva, crops the mask's bounding box and
regresses Hb from the crop.}
\label{fig:pipeline}
\end{figure*}
% Condensed by Claude (23 Nov) from the existing text (AI-written); "pulling
% the eyelid down" is Anubhav's phrase from Section I. The JPEG sentence is
% from Biraj's revised VII-A (23 Nov), verbatim.
Fig.~\ref{fig:pipeline} shows the pipeline. A photograph of the lower eyelid,
pulled down, is taken in the application and sent to the server. An Attention
U-Net segments the conjunctiva, and the mask's bounding box with a 20-pixel
margin is cut from the full-resolution photograph. It is then JPEG-encoded in
memory at quality 90, so that it matches the compression of the training crops.
A DenseNet-121 estimates Hb from this crop, and the
application shows the estimate with a screening flag. Sections~V--VI are an
independent study of conjunctiva extraction on the public datasets; the deployed
segmenter is the architecture selected there, an Attention U-Net with an
EfficientNet-B0 encoder, trained separately on the hospital masks
(Sections~VII--VIII).

% ============================================================ V. EXTRACTION (part: extraction)
\section{Conjunctiva Extraction}
\subsection{Quality Control}
% Written by Anubhav (Q&A, 22 Nov). Edits: grammar, tense, LaTeX. Facts added
% by Claude: the public-photo range for colour cast (measured on all 764
% public photographs with preprocess.cast_ratio), what the coloured-light range
% was measured on (experiments/04), and "before any colour is measured".
Because anaemia detection depends on colour, it becomes very important to
check the quality of the image taken. A blurred picture or a picture clicked in
a coloured light can give a wrong reading, which would be worse than getting no
reading. So, to prevent wrong readings, each picture goes through four checks,
namely:
\begin{itemize}
\item \emph{Focus}, which uses the Laplacian, an edge detector. A sharp picture
has strong edges, leading to high variance, as its Laplacian values are spread
widely. On the other hand, a blurred picture has weak edges and low variance.
We reject images with variance below 35.
\item \emph{Mask size}, where we check the percentage of the image the mask
covers. If it covers almost nothing or almost everything, there is a high
chance something is wrong. This can happen in cases where the eyelid was not
pulled down or there is no eye in the picture. We reject images with mask size
below 0.2\% or above 85\% of the frame.
\item \emph{Clipping}, where a flash or a very bright light can turn pixel
values in the image to pure white (all channels 255). More than 5\% of the frame
is where we keep our limit. This is just a generous guess, as we could not find
such a case to calibrate it.
\item \emph{Colour cast}. In this, we divide the mean of the brightest colour
channel by the mean of the dimmest. Normal photos usually score 1.02--2.04 (all 764 public photographs).
Pictures under a strongly coloured light can score 22.6--206 (one volunteer's
eye under five coloured lights). We use a limit of 3 to separate the two.
\end{itemize}
An image that fails any check is rejected, and the check it failed is returned,
before any colour is measured.

\subsection{Segmentation}
% Written by Anubhav (Q&A, 22 Nov); edits: I->we, lists inline, citations, typos.
We compared two types of extractor: rule-based and learned. The rule-based
extractor finds the conjunctiva using rules, limits and thresholds on the
colour of each pixel, refined with seeded GrabCut~\cite{rother2004grabcut}.
Where it fails is that in anaemia the conjunctiva itself is
very pale, and if we define rules to find such pale tissue, the extractor might
leak into the sclera and skin. Because colour is the most important signal for
estimating the Hb level, we assume that precision matters more than recall, so
rule-based extractors were deemed not ideal for this task.

For the learned extractor, we tested six decoders:
U-Net~\cite{ronneberger2015}, UNet++~\cite{zhou2018unetpp}, Attention
U-Net~\cite{oktay2018}, LinkNet~\cite{chaurasia2017}, FPN~\cite{lin2017fpn} and
PSPNet~\cite{zhao2017pspnet}. All six used the same ImageNet-pretrained
EfficientNet-B0 encoder~\cite{tan2019efficientnet}, so only the decoder differed
between them. For training, we used the publicly available
Eyes-Defy-Anemia~\cite{dimauro2023} and Mendeley~\cite{rahman_mendeley}
datasets. From Eyes-Defy-Anemia we used the 101 traced masks (Section~III); for
Mendeley, where every image has two annotators, we used the intersection of
their masks, the region both marked as conjunctiva. Each model was trained for
40 epochs at $320 \times 320$ pixels with AdamW (learning rate $3 \times
10^{-4}$, cosine schedule, batch size 8) on binary cross-entropy plus Dice loss,
with random crops, flips, rotations of up to $45^\circ$ and colour jitter, and
the epoch with the best validation Dice was kept.
% Facts filled in by Claude (22 Nov, at Anubhav's request): shared-encoder
% sentence, masks sentence, training-settings sentence. Sources:
% scripts/train_segmenter.py defaults, experiments/24 run.py (rotation
% augmentation, 40 epochs), scripts/prepare_conjunctiva_seg.py (intersection).

\subsection{Evaluation Protocol}
% Written by Anubhav (Q&A, 22 Nov). Edits: grammar, LaTeX, equation references.
% Facts added by Claude: "only", "pixel", the buffer sentence, SD (not
% variance), the statistical-test sentence, the fixed-weight sentence.
To measure the correctness of a mask, we use $P$, the predicted mask, and $T$,
the annotated mask. The mask is then scored with metrics like Dice and IoU (1).
These both count only the conjunctiva and rank models in the same way, but IoU
is stricter. Precision (2) is the share of the predicted mask that is actually
the conjunctiva. Low precision means there is a spill. Recall is the share of
the annotated conjunctiva that was found. Low recall means the mask is too
small. HD95 is the 95th-percentile distance between the two outlines in pixels.
It tells us whether the outline is in the right place. The 95th percentile is
used rather than the maximum so that a few outlier pixels do not dominate the
result. We are not using pixel accuracy, because the conjunctiva covers only
about 1.4\% of the frame, so predicting everything as background already gives
98.6\% accuracy.
\begin{equation}
\mathrm{Dice} = \frac{2\,|P \cap T|}{|P| + |T|}, \qquad
\mathrm{IoU} = \frac{|P \cap T|}{|P \cup T|}
\end{equation}
\begin{equation}
\mathrm{Precision} = \frac{|P \cap T|}{|P|}, \qquad
\mathrm{Recall} = \frac{|P \cap T|}{|T|}
\end{equation}
\begin{equation}
F_{0.5} = \frac{1.25\,|P \cap T|}{1.25\,|P \cap T| + 0.25\,|T \setminus P| + |P \setminus T|}
\end{equation}

The Mendeley dataset has many images of the same person's eyes, taken from
different angles, so while making the split, we needed to make sure the same
person's images remain together, either in training, validation or testing.
This is because if the split causes the same person's eye to be present in both
training and testing, this could cause the model to score higher. To prevent
this, Mendeley is dealt out in blocks of 24 consecutive files, because
neighbouring groups are often the same person's other eye. Files within four
numbers of a test image are also removed from training. Each dataset gives
20\% of its images to testing, 20\% to validation to find the optimal number of
epochs, and the rest to training. This is repeated over five splits, and the
mean and standard deviation of the scores are reported. Differences are tested
with the corrected resampled $t$-test~\cite{nadeau2003}, because the splits
share most of their training images, with Holm's correction~\cite{holm1979}
when several comparisons are made together.

As stated in Section~V-B, we assume precision matters more than recall, so we
also evaluated the models with $F_{0.5}$ (3)~\cite{vanrijsbergen1979}.
$F_{0.5}$ gives more weight to precision compared to recall, twice as much. The
weight was fixed before any model was scored. Each model's threshold is chosen
on its own validation images, because one model never gave any pixel a
probability above 0.53.

\subsection{Colour Correction}
% Written by Anubhav (Q&A, 22 Nov). Edits: grammar, method names and citations,
% "proposed" -> "compared" (gray-world and white-patch are existing methods).
% Facts added by Claude: where the sclera reference is taken, "so the
% segmenter's input and masks are unchanged", the gains sentence, the figure
% reference, and why the plastic is left out of the white patch.
In this section, we talk about the colour correction methods we compared. The
colour correction part is kept in the preprocessing of the image because it
requires the whole image. The error in colour of the image is detected from the
surroundings, like the entire frame or the eye. This cannot be done after
cropping, because in that case everything else is marked as background and
removed. The methods compared are:
\begin{itemize}
\item \emph{Gray-world}~\cite{buchsbaum1980}, where we assume the whole picture
averages to grey.
\item \emph{White-patch}~\cite{land1971} assumes the brightest pixels in the
picture are white. The plain version (max-RGB) changes nothing here, because in
96\% of photos every channel reaches 255 somewhere (clipped). Ours takes the
brightest 1\% of unclipped pixels (any channel at 250 or above is skipped), and
only above the lower lid, a region found using the mask.
\item The \emph{sclera} method assumes the sclera of the eye is white; we take
its brightest part, in a band above the conjunctiva.
\end{itemize}
This is applied to the entire picture after the segmentation, so the
segmenter's input and masks are unchanged. Each method gives one gain per colour
channel, chosen to keep the overall brightness.

To check the correctness of the colour correction methods, we used the
Eyes-Defy-Anemia pictures, which were taken with a device that has a white
plastic edge below the eye (visible along the bottom of each photograph in
Fig.~\ref{fig:light}). The camera's white balance was fixed and no flash was
used, so the plastic's colour records the colour of the light only. The angle
between the estimated light colour and the plastic's colour was measured to find
the error. The angle ignores brightness, and a smaller angle is better. The
white patch leaves out the plastic, since about half of the brightest pixels of
the whole frame lie on it.

\section{Extraction Results}
\subsection{Rule-Based and Learned Extraction}
\begin{table}[t]
\caption{Rule-Based and Learned Extraction against Annotated Masks
(Dice, Mean over Five Person-Safe Splits)}
\label{tab:rules}
\centering
\footnotesize
\begin{tabular}{@{}lcc@{}}
\toprule
Extractor & Close-ups & Whole-face \\
\midrule
Rule-based (seeded GrabCut) & 0.476 $\pm$ 0.022 & 0.034 $\pm$ 0.021 \\
U-Net & \textbf{0.832} & 0.825 \\
Attention U-Net & \textbf{0.832} & \textbf{0.831} \\
\bottomrule
\end{tabular}
\end{table}
\begin{figure}[t]
\centering
\fig{0.85\columnwidth}{extractor_comparison.png}
\caption{Rule-based (top) and learned U-Net (bottom) masks on three
Eyes-Defy-Anemia photographs drawn at random from those never used in
training, because their annotations were rectangles (CC~BY-SA~4.0).}
\label{fig:extractors}
\end{figure}
% Written by Anubhav (Q&A, 22 Nov). Edits: grammar, table/figure references.
% Facts added by Claude: "on all 100 test cases", "about" 0.83 (U-Net 0.825,
% Attention U-Net 0.831; experiment 27), the closing caveat sentence.
The images used to test both the extractors are the same
(Table~\ref{tab:rules}). In the case of learned extractors, we test the models
trained on the training split only. The results reported clearly showed the
learned extractors to be much better than the rule-based ones. When tested on
close-ups in the Eyes-Defy-Anemia dataset, the rules reach a Dice of 0.48
against 0.83 for the networks (U-Net and Attention U-Net both). The network
wins across all five splits, on all 100 test cases. We then noted the
precision: the rule-based masks show around 0.46, meaning more than half the
mask is actually not the conjunctiva. As discussed, spill is worse than
missing. On the full-face photographs in Mendeley, the rules fail at 0.03. Their
thresholds assume the picture to be mostly eye, so they grab skin on the face
or a finger instead. The networks score about 0.83 there too.
Fig.~\ref{fig:extractors} shows an example of masks on a few pictures. The rules
do remain steadier when a photo is rotated or zoomed, which this comparison does
not test.

\subsection{Architecture Comparison}
\begin{table}[t]
\caption{Segmentation on People Not Seen in Training}
\label{tab:arch}
\centering
\footnotesize
\setlength{\tabcolsep}{3.5pt}
\begin{tabular}{@{}lcccccc@{}}
\toprule
Architecture & Dice & IoU & Prec. & Rec. & HD95 & Rank \\
\midrule
UNet++ & 0.833 $\pm$ 0.017 & 0.722 & 0.820 & 0.865 & 3.3 & 1.4 \\
Attention U-Net & 0.831 $\pm$ 0.012 & 0.719 & 0.824 & 0.854 & 3.5 & 2.2 \\
U-Net & 0.826 $\pm$ 0.018 & 0.712 & 0.825 & 0.844 & 3.4 & 2.8 \\
FPN & 0.821 $\pm$ 0.018 & 0.704 & 0.806 & 0.855 & 3.4 & 4.0 \\
LinkNet & 0.810 $\pm$ 0.021 & 0.689 & 0.767 & 0.880 & 3.6 & 4.6 \\
PSPNet & 0.764 $\pm$ 0.029 & 0.631 & 0.730 & 0.829 & 5.5 & 6.0 \\
\bottomrule
\end{tabular}

\vspace{2pt}
\parbox{\columnwidth}{\scriptsize Mean over five person-safe splits (118--140
test images each), Dice with SD across splits. HD95 in pixels on the 320-pixel
grid; Rank is the mean rank by Dice.}
\end{table}
% Written by Anubhav (Q&A, 22 Nov). Edits: grammar, "U-Net++" -> "UNet++" (as in
% the table), citation style, "small" -> "not statistically significant" for the
% p-values. Facts added by Claude: "FPN and LinkNet follow, and PSPNet is last in
% every split"; recall 0.880 / precision 0.767; Holm-corrected; data and training
% length as further differences (Dhalla: DenseNet-121, 200 epochs, 324 photos of
% 162 children; ours: EfficientNet-B0, 40 epochs); SD 0.012 vs 0.017; F0.5 at
% the default threshold (0.825). Source: experiment 24.
We compared six architectures: U-Net, UNet++, Attention U-Net, FPN, LinkNet and
PSPNet. They were trained and tested on all five person-safe splits. The mean
and SD of Dice across the splits are shown in Table~\ref{tab:arch}.

As recorded, the top three are within 0.007 Dice of each other, which is too
small a range to decide a winner. FPN and LinkNet follow, and PSPNet is last in
every split. Dhalla et~al.~\cite{dhalla2023} chose LinkNet, but in our setting
it never came first in any of the splits. Even though LinkNet gives the highest
recall (0.880), it gives lower precision (0.767), meaning that even though it
grabs the conjunctiva, it also spills, which is the worse case for us. On the
person-safe splits, its gap to the U-Net family is not statistically
significant (Holm-corrected $p$ 0.18--0.43), and their data, encoder and
training length differ from ours, so we say their result was `not reproduced'
instead of calling it wrong.

We decided to go with Attention U-Net because it ties for the best Dice, has the
smallest spread across splits (SD 0.012 against 0.017 for UNet++) and has the
highest $F_{0.5}$ at the default threshold (0.825).

\subsection{Checks on the Evaluation}
% Paragraphs 2-5 written by Anubhav (Q&A, 22 Nov). Edits: grammar, tense,
% contractions, "color" -> "colour", "When Mendeley was split randomly" -> "When
% the images were split randomly" (the 129 test images per split come from both
% datasets). Added by Claude: the opening sentence, the italic labels, "of the
% 129 ... in each split", "Holm-corrected", "on its validation images", the
% test-F0.5 clause (29 of 30 models; experiment 25).
% The last paragraph was written by Claude at the author's request, in the
% author's plain style. Sources: experiment 22 (grid bound: 0.992 Mendeley,
% 0.998 Eyes-Defy-Anemia); the 16 Sep annotator analysis (547 pairs: second
% annotator larger in 99%, 98.1% of the smaller inside the larger).
We ran four checks on our own evaluation, and each one changed a number we
would otherwise have reported.

\emph{One split against five.} On a single split, Attention U-Net beats U-Net
by 0.0105 Dice, which is significant ($p = 0.007$). When we repeated the
process over five splits, this advantage was reduced to 0.0019 ($p = 0.66$).
That split was simply the lucky split, which can make a difference look real
when it is not.

\emph{Leakage.} When the images were split randomly, 83--93 of the 129 test
images in each split were of a person whose other images were in the training
set as well. The model had effectively seen the eye before. Person-safe
splitting lowered Dice by only 0.001--0.018 and did not change the order. The
difference it made was that LinkNet's losses were significant before
(Holm-corrected $p$ 0.03--0.05) but not after.

\emph{Train/serve mismatch.} The segmenter was trained on raw pictures, but
when fed colour-corrected images, its Dice fell by 0.012. So the segmenter now
operates on raw photos, and colour correction is done later.

\emph{Threshold.} The network gives each pixel a probability, and 0.5 is the
usual cutoff. Choosing each model's cutoff to maximise $F_{0.5}$ on its
validation images moved it to about 0.97 and raised $F_{0.5}$ on the test
images in all 29 models where it moved. Maximising $F_{0.5}$ caused precision
to move up by 0.06--0.12, but cost 0.02--0.03 Dice. Choosing the cutoff for
Dice gained nothing, so 0.5 is fine as a threshold when Dice is the goal.

We also checked two things that affect how the scores should be read. First,
the networks predict masks at $320 \times 320$ pixels, while the annotations
are drawn on the full-size photograph. Even a perfect mask at $320 \times 320$
would still score about 0.99 Dice against the full-size annotation (0.992 on
Mendeley, 0.998 on Eyes-Defy-Anemia), so the smaller size costs very little
compared to the 0.83 the networks reach. Second, the two Mendeley annotators
agree with each other at only 0.736 Dice, which looks low. But their masks are
nested: the second annotator's mask is the larger one in 99\% of the images,
and 98\% of the smaller mask lies inside the larger one. So they agree on where
the conjunctiva is, and only differ on how far its upper edge goes, since the
conjunctiva has no sharp upper edge and fades into the fold behind the lid.
This also means the intersection we train on (Section~V-B) is almost the same
as the smaller of the two masks.

\subsection{Estimating the Light}
\begin{table}[t]
\caption{Median Angle between Each Estimate of the Light and the Device
Plastic}
\label{tab:light}
\centering
\footnotesize
\begin{tabular}{@{}lcc@{}}
\toprule
Estimate & India (90) & Italy (120) \\
\midrule
None & 4.7$^\circ$ & 2.8$^\circ$ \\
Gray-world & 4.9$^\circ$ & 5.2$^\circ$ \\
White-patch, above the lower lid & \textbf{1.5$^\circ$} & \textbf{0.8$^\circ$} \\
Sclera & 2.3$^\circ$ & 0.9$^\circ$ \\
\bottomrule
\end{tabular}
\end{table}
\begin{figure*}[t]
\centering
\fig{0.75\textwidth}{white_balance_comparison_india_report.png}
\caption{One Indian Eyes-Defy-Anemia photograph under each correction, with the
same mask in all four. ``Plastic off grey'' is the device plastic's angle from
neutral after correction (CC~BY-SA~4.0).}
\label{fig:light}
\end{figure*}
% Written by Anubhav (Q&A, 22 Nov). Edits: grammar, "Grey-World" -> "Gray-world"
% (as in Table IV and V-D), "averages out to Grey-World" -> "averages out to
% grey", "a close-up of face or eye" -> "a close-up of the eye" (Table IV is
% Eyes-Defy-Anemia close-ups only). Added by Claude: what the two arrows are,
% "whatever its brightness", "(India/Italy)", the figure reference, the sentences
% after "yellower with age" up to and including the default sentence, and
% "(Section X)". Source: experiment 26 (sclera b* minus plastic b*, per decade,
% Hb and sex held: +0.64 India, +0.62 Italy; sclera a* against Hb rho +0.51 /
% +0.44; India, brightest 1% above the lid: older patients +0.47 g/dL per decade).
To measure how close or far we are, we represent colours as vectors in 3-D
space, with R, G and B as its axes. The score is measured by the angle between
the two arrows, the estimated light and the plastic, and 0$^\circ$ represents
exactly the same colour, whatever its brightness. Table~\ref{tab:light} gives
the median angle on 90 photos from India and 120 from Italy (Eyes-Defy-Anemia).

In the case of no correction, the recorded angle was 4.7$^\circ$/2.8$^\circ$
(India/Italy). In the case of Gray-world, we recorded 4.9$^\circ$/5.2$^\circ$,
worse than doing nothing. A possible explanation is that Gray-world assumes the
photo averages out to grey, but a close-up of the eye would mostly have pink
tissue and skin. The white patch gives 1.5$^\circ$/0.8$^\circ$, the best, and
the sclera gives 2.3$^\circ$/0.9$^\circ$, the second best (Fig.~\ref{fig:light}).

The sclera is not really perfectly white, though. It becomes yellower with age:
measured against the plastic, its yellowness ($b^*$ in CIELAB) rises by 0.64
(India) and 0.62 (Italy) per decade, comparing people with the same Hb and sex.
Its redness ($a^*$) also follows Hb, so correcting by it removes some of the Hb
signal too. Since the brightest pixels above the lid lie mostly on the sclera,
the white patch inherits part of this: in India it would make older patients
read lower by about 0.47~g/dL per decade. We still made the white patch the
extractor's default, because it estimates the light most closely.

We have not yet measured the effect of these corrections on the haemoglobin
prediction. It is included in the future plans for this project (Section~X).

% ============================================================ VII. HAEMOGLOBIN ESTIMATION (part: Hb model)
\section{Haemoglobin Estimation}
% VII-VIII: Biraj's text, his revised version of 23 Nov. Merged by Claude on 23 Nov
% by deletion only, plus one "and": the patient split (in III), pixel specificity
% and accuracy (V-C explains why accuracy is not used), the JPEG step (in IV and
% VIII-A), target-standardisation constants, seed IDs, other models' MAEs, repeated
% numbers, the predicted-vs-measured and residual plots (same as Fig. 6 a and b),
% and server internals. Figure widths set to fit 8 pages.
\subsection{Hospital Segmentation and Crops}
All 288 hospital photographs were outlined by hand in an annotation tool we
built for this purpose. Every mask starts blank and no model prediction is
shown, so the annotator is not influenced by a model. Eighty-nine portrait
photographs store their pixels rotated with an EXIF orientation flag, so we
apply EXIF transposition when loading each image, before pairing it with its
mask.
Based on the architecture comparison in Sections~V--VI, we trained an Attention
U-Net with an ImageNet-pretrained EfficientNet-B0 encoder on these clinical
photographs at $320 \times 320$ resolution. The loss was binary cross-entropy
plus soft Dice, and we used AdamW, batch size 2, horizontal flips, $\pm
8^\circ$ rotation and $\pm 10\%$ brightness changes. Training was stopped early
based on validation Dice (best epoch 18, 0.8301; Fig.~\ref{fig:training_curves}).
On the 43 test patients, who were never seen during training, the clinical
Attention U-Net reached a Dice of $0.8273 \pm 0.0694$, IoU $0.7112 \pm 0.0962$,
precision $0.8180 \pm 0.1039$ and recall $0.8555 \pm 0.1123$.
The crop given to the regressor is the bounding box of the predicted mask with a
20-pixel margin on all four sides, cut from the full-resolution photograph.

\begin{figure}[t]
\centering
\fig{0.9\linewidth}{training_curves.png}
\caption{Training and validation loss and Dice coefficient across epochs for the selected segmentation model.}
\label{fig:training_curves}
\end{figure}

\subsection{Regression Model}
Each crop is resized to $224 \times 224$ and normalised with ImageNet
statistics. No colour correction is applied. We compared six ImageNet-pretrained
backbones: ResNet-18~\cite{he2016resnet}, VGG-16~\cite{simonyan2015vgg},
DenseNet-121~\cite{huang2017densenet}, EfficientNet-B0~\cite{tan2019efficientnet},
MobileNetV3-Large~\cite{howard2019mobilenetv3} and
ConvNeXt-Tiny~\cite{liu2022convnext}. Each has a single-output linear regression
head. The backbone is kept frozen for the first 10 epochs, and after that its
final block is fine-tuned (learning rate $10^{-4}$ for the backbone and
$10^{-3}$ for the head, AdamW, weight decay $10^{-4}$, batch size 8, Smooth L1
loss, up to 40 epochs, with early stopping after 7 epochs without improvement in
validation error). All BatchNorm statistics stay frozen, with the layers in
evaluation mode, throughout training. The targets are standardised using the
mean and standard deviation of the training split.
Data augmentation consists of $\pm 5^\circ$ rotation, horizontal flips, and
$\pm 5\%$ brightness and contrast changes.

\subsection{Training and Evaluation}
The 279 patients with a valid clinical Hb record were split 195/43/41 (training,
validation, test), which matches the patient-wise segmentation split exactly. We
compared the architectures on validation mean absolute error (MAE). The three
best candidates (DenseNet-121, MobileNetV3-Large and VGG-16) were then run with
three random seeds. DenseNet-121 had the lowest
mean validation MAE at every seed ($0.8760 \pm 0.0895$~g/dL). We
selected the checkpoint with the lowest validation MAE, which was epoch 9
(validation MAE 0.7736~g/dL). The
held-out test split was not used at all for choosing the architecture or tuning
hyperparameters, and the final model was evaluated on it only once. The naive
baseline predicts the training mean (11.3764~g/dL) for every patient. We report
MAE, median absolute error (MedAE), RMSE, $R^2$, Pearson correlation $r$, the
maximum absolute error, and Bland--Altman 95\% limits of agreement~\cite{bland1986}.

\subsection{Results}
\begin{table}[t]
\caption{Haemoglobin Estimation Performance on Hospital Dataset}
\label{tab:hb}
\centering
\footnotesize
\begin{tabular}{@{}lcccc@{}}
\toprule
Model & MAE & MedAE & RMSE & $R^2$ \\
\midrule
\multicolumn{5}{@{}l}{\emph{Validation split, 43 patients (first seed 20261116)}} \\
DenseNet-121 & \textbf{0.939} & 0.757 & 1.337 & 0.282 \\
MobileNetV3-Large & 0.953 & 0.777 & 1.323 & 0.297 \\
EfficientNet-B0 & 0.961 & 0.814 & 1.341 & 0.278 \\
VGG-16 & 0.992 & 0.742 & \textbf{1.270} & \textbf{0.352} \\
ResNet-18 & 1.028 & 0.835 & 1.317 & 0.304 \\
ConvNeXt-Tiny & 1.065 & 0.841 & 1.399 & 0.215 \\
\midrule
\multicolumn{5}{@{}l}{\emph{Held-out test split, 41 patients (final evaluation)}} \\
Training mean baseline & 1.371 & \textbf{0.676} & 1.915 & $-0.002$ \\
DenseNet-121 (epoch 9) & \textbf{1.345} & 1.066 & \textbf{1.701} & \textbf{0.210} \\
\bottomrule
\end{tabular}
\end{table}
\begin{figure}[t]
\centering
\fig{0.65\columnwidth}{absolute_error_distribution.png}
\caption{Distribution of absolute haemoglobin prediction errors on the held-out test set of 41 patients.}
\label{fig:abs_error_dist}
\end{figure}

\begin{figure}[t]
\centering
\fig{\columnwidth}{bland_altman.png}
\caption{DenseNet-121 on the 41 held-out test patients. (a) Predicted against
measured Hb; dotted lines mark the 11.0~g/dL screening threshold. (b) Bland--Altman
difference plot: mean bias of $+0.408$~g/dL (solid) and 95\% limits of agreement
from $-2.868$ to $+3.684$~g/dL (dashed).}
\label{fig:ba}
\end{figure}
Table~\ref{tab:hb} gives the results. On the 41 unseen test patients, it gave an MAE of 1.3452~g/dL, an
RMSE of 1.7005~g/dL and an $R^2$ of 0.2097, with a statistically significant
positive Pearson correlation ($r = 0.4705$, $p = 0.0019$).

The drop in MAE is small ($-0.026$~g/dL), but the drop in RMSE is larger
($-0.214$~g/dL). This suggests that the network mainly helps by limiting the
very large errors: the maximum error went down from 5.276~g/dL for the baseline
to 4.500~g/dL for DenseNet-121. The baseline actually has a lower median
absolute error (0.676~g/dL against 1.066~g/dL), because the Hb values in the
test cohort are mostly concentrated around 10.7--12.0~g/dL. However, the
baseline cannot tell patients apart at all ($R^2 < 0$). The image-based
predictions, on the other hand, correlate positively with laboratory Hb, which suggests that the conjunctival crop does carry
some optical information related to haemoglobin concentration.

The predicted-versus-measured plot (Fig.~\ref{fig:ba}a) shows that the
predictions are fairly scattered and are pulled toward the training mean. They
range from 9.42 to 13.45~g/dL, while the laboratory values span 6.10 to
15.00~g/dL. The distribution of absolute errors
on the test set is in Fig.~\ref{fig:abs_error_dist}. The Bland--Altman plot
(Fig.~\ref{fig:ba}b) shows a positive mean bias of $+0.4078$~g/dL, with 95\% limits of agreement from $-2.8682$ to
$+3.6837$~g/dL. At the WHO 11.0~g/dL screening threshold, 7 of the 15 anaemic
test patients were correctly flagged and 22 of the 26 non-anaemic patients were
classified as normal. The validation MAE (0.774~g/dL) should not be mistaken for
how well the model generalises. The test MAE of 1.345~g/dL is the more honest
number, and it shows that the system is at present a preliminary, non-invasive
screening prototype and not a replacement for venous blood analysis.

\section{Mobile Application}
\subsection{Architecture}
The mobile application is built in Flutter (supporting Android, iOS and Web) and
talks to a production FastAPI backend that runs on a local or cloud server.
When an image is uploaded, the server applies EXIF transposition, segments the
conjunctiva, extracts the crop with the 20-pixel margin, compresses it in memory
as JPEG at quality 90, and runs the DenseNet regression. It then returns a JSON
response containing the predicted Hb (g/dL) and previews of both the isolated crop
and the segmentation overlay.

\subsection{Screening Output}
The user interface (Fig.~\ref{fig:app}) gives step-by-step photo instructions
(everting the eyelid, using ambient daylight, checking the focus) and shows the
resulting Hb estimate along with previews of the tissue. The system keeps the Hb
prediction separate from the risk assessment. First, the DenseNet-121 model uses
no demographic inputs, and its Hb estimate (g/dL) is based only on the cropped
RGB image of the conjunctiva. Second, the screening threshold is chosen on the
client side from the demographics the patient reports, following WHO guidelines:
12.0~g/dL for non-pregnant adult women (as shown for the 37-year-old female in
Fig.~\ref{fig:app}), 13.0~g/dL for adult men, and 11.0~g/dL for pregnant women
or when the demographic fields are left empty. Third, the screening flag is
produced by comparing the predicted Hb directly with the chosen threshold
($\widehat{\text{Hb}} < \text{Threshold}$ means there is a risk of anaemia). This
tells the user to get a confirmatory clinical test.
\begin{figure}[t]
\centering
\fig{0.45\linewidth}{eyecheck_application.png}
\caption{EyeCheck mobile application interface showing the photo input, haemoglobin screening result, segmentation mask, and extracted conjunctiva region.}
\label{fig:app}
\end{figure}

\section{Discussion and Limitations}
% Assembled by Claude (23 Nov). Anubhav's own sentences from VI-D and V-A:
% "We have not yet measured the effect of these corrections on the haemoglobin
% prediction", "is not really perfectly white ... becomes yellower with age",
% "just a generous guess, as we could not find such a case to calibrate it".
% The rest is condensed from the earlier draft and Section VII (AI-written);
% "pulled toward" and "7 of the 15" follow Biraj's revised VII-D (23 Nov).
Our extraction and colour-correction results come from public photographs taken
with other phones, and have not yet been repeated on the hospital phones. We
have not yet measured the effect of these corrections on the haemoglobin
prediction: the regressor is trained on uncorrected crops, and the sclera, on
which the white patch partly lands, is not really perfectly white, as it
becomes yellower with age. The clipping limit in quality control is just a
generous guess, as we could not find such a case to calibrate it. The crop is a
padded bounding box, so it contains lid skin and sclera. The regressor's
predictions are pulled toward the training mean, and at 11.0~g/dL only 7 of the
15 anaemic test patients were flagged. All datasets are small, and nothing here
is clinically validated.

% ============================================================ X. CONCLUSION (shared)
\section{Conclusion and Future Work}
% Assembled by Claude (23 Nov) from Anubhav's own sentences in VI-A to VI-D
% ("clearly showed the learned extractors to be much better than the rule-based
% ones", "too small a range to decide a winner", "can make a difference look
% real when it is not", "the white patch gives ... the best"). The hospital, Hb
% and future-work sentences are condensed from the earlier draft (AI-written).
% Encoder sentence (23 Nov): Anubhav's idea ("using different encoders for
% Networks"); the "since all six decoders..." clause is Claude's, from V-B.
The results clearly showed the learned extractors to be much better than the
rule-based ones (Dice 0.83 against 0.48 on close-ups). The top three
architectures are within 0.007 Dice of each other, which is too small a range to
decide a winner, and a single split can make a difference look real when it is
not. Of the colour corrections, the white patch gives the best estimate of the
light. The segmenter selected on the basis of this study reaches a Dice of 0.83
on hospital photographs, and a DenseNet-121 estimates Hb with an MAE of
1.35~g/dL, slightly better than predicting the mean. Next, we will retrain the
Hb regressor on white-patch-corrected crops, repeat the extraction study on the
hospital photographs, and evaluate on an independent cohort. We also plan to
try different encoders for the networks, since all six decoders in our study
shared the same EfficientNet-B0 encoder.

\section*{Acknowledgment}
We thank our supervisor, Amit Chintamani Awekar, for guidance throughout the
project; Dr.~Sandip Paul, MBBS, MMST (IIT KGP), IMO (Grade~II) at ESIC Medical
College and Hospital (ESICMCH), Beltola, Guwahati, for clinical facilitation, data
acquisition, and medical oversight;
and the authors of the Eyes-Defy-Anemia and Mendeley datasets for making them
public.

\bibliographystyle{IEEEtran}
\bibliography{refs}

\end{document}
